\subsection{An order-nine class character and three exact balanced candidates}

\begin{statement}{A.174}{Proved}
{The level-three balanced target reduces to three ordinary class tests}
Retain the actual fields, primes and completed groups of A.165--A.170.
All ideal classes in this statement mean their three-primary projections.
In \(A=\operatorname{Cl}(K_3)\{3\}\), put
\[
 \begin{gathered}
 P=[P_{2,0,3}],\qquad L_q=[P_{q,0,3}],\qquad Q_q=[P_{q,2,3}],\\
 R_q=Q_q-2L_q\quad(q=11,59),\qquad
 G=\langle P,R_{11},R_{59}\rangle,\qquad
 \mathcal R=\langle R_{11},R_{59}\rangle .
 \end{gathered}
\tag{A.174.1}
\]
There is an actual surjective homomorphism
\[
 \theta:G\longrightarrow\Z/9,\qquad
 \theta(P)=-3,\quad\theta(R_{11})=-4,\quad\theta(R_{59})=1 .
\tag{A.174.2}
\]
Its reduction modulo three is exactly the restriction of the
unramified genus character \(\chi_3\) of A.170.
Both \(R_{11}\) and \(R_{59}\) have exact order nine.
For
\[
              r=R_{11}-5R_{59},\qquad e'=P+3R_{59},
\tag{A.174.3}
\]
one has \(3r=0\), and the following decompositions and quotient map
retain the actual ideal classes:
\[
 \begin{gathered}
 \mathcal R=\langle R_{59}\rangle\oplus\langle r\rangle,
       \qquad \langle R_{59}\rangle\simeq\Z/9,\\
 G=\langle R_{59}\rangle\oplus\langle e',r\rangle,\\
 \operatorname{im}(G\longrightarrow C_3^{\rm bal})
       \simeq \langle e',r\rangle/\langle r,3e'\rangle .
 \end{gathered}
\tag{A.174.4}
\]
In particular this last group is cyclic of order at most three,
with the balanced target \(e_3\) as its specified generator.

Define the three literal fractional ideals
\[
 \begin{gathered}
 \mathfrak R_q=P_{q,2,3}P_{q,0,3}^{-2},\\
 \mathfrak I_j
   =P_{2,0,3}\,\mathfrak R_{11}^{\,j}
                    \mathfrak R_{59}^{\,3-5j},
                  \qquad j=0,1,2 .
 \end{gathered}
\tag{A.174.5}
\]
Each has rational ideal norm two.  The exact reduction is
\[
 e_3=0
 \quad\Longleftrightarrow\quad
 [\mathfrak I_j]_{(3)}=0\text{ for some }j\in\{0,1,2\}.
\tag{A.174.6}
\]
Equivalently, one of these three fractional ideals has a principal
positive integral power whose exponent is prime to three.
The result supplies neither a vanishing nor a nonvanishing assertion
for the three classes in A.174.6.
\end{statement}

\begin{proof}
\emph{The homogeneous form of the actual norm and residue maps.}
Use the elements \(\beta_1,\beta_2,U,h_{11},h_{59},\eta\) of A.165.
In particular
\[
 N_{2,1}\beta_2=\beta_1,\qquad
 N_{2,1}\eta=Uh_{11}h_{59},\qquad
 (\vartheta_{11}(\eta),\vartheta_{59}(\eta))=(0,2).
\]
Here \(\vartheta_q\) denotes the actual stable tame norm character
at the quadratic \(q\)-place, with the chosen residue generator
of A.165.  Its other values needed below are
\[
 \vartheta(\beta_2)=(2,1),\quad
 \vartheta(h_{11})=(1,0),\quad
 \vartheta(h_{59})=(0,1).
\tag{A.174.7}
\]
These are additive characters with values in \(\F_3\).
The elements \(h_q\) in this display are the scalar images from
\(K_1\); their relative norm to \(K_1\) is \(h_q^3\).

Suppose a completed \(S\)-unit \(x\) of \(K_2\) has relative norm
\[
           \beta_1^k U^z h_{11}^{\,b}h_{59}^{\,c},
                       \qquad k,z,b,c\in\Z_3.
\tag{A.174.8}
\]
Multiplying by \(\beta_2^{-k}\) puts this in precisely A.165's
homogeneous norm fibre.  Its exact first norm image gives
\(b-z,c-z\in3\Z_3\).  An explicit element of that fibre is
\[
 x_0=\beta_2^k\eta^z
        h_{11}^{(b-z)/3}h_{59}^{(c-z)/3}.
\tag{A.174.9}
\]
Every other element differs by a completed relative norm-one
\(S\)-unit, whose two characters vanish by A.159.
Consequently the two characters of \(x\) are literally
\[
      2k+(b-z)/3,\qquad k+2z+(c-z)/3\quad\text{in }\F_3.
\tag{A.174.10}
\]
If \(x\) is a norm from \(K_3\), these are zero.
Multiplying the two equalities by three gives
\[
       b\equiv z+3k,\qquad c\equiv4z+6k\pmod9,
       \qquad c-4b-3k\equiv0\pmod9.
\tag{A.174.11}
\]
This is a necessary local condition for the displayed actual norm;
the argument does not assert global norm existence from this condition.

\emph{Descent of the order-nine character through actual principal ideals.}
Let \(k,b,c\) be integers with
\(kP+bR_{11}+cR_{59}=0\) in \(A\).
The corresponding fractional ideal has prime-to-three class order.
Choose a positive integer \(h\), prime to three, and
\(w\in K_3^\times\) with
\[
 \operatorname{div}(w)
 =h\bigl(kP_{2,0,3}
       +b(P_{11,2,3}-2P_{11,0,3})
       +c(P_{59,2,3}-2P_{59,0,3})\bigr).
\tag{A.174.12}
\]
This is an actual \(S\)-unit, also for negative exponents.
The rational norm has absolute value \(2^{hk}\).
All three indicated prime branches have residue-degree-one
ideal norm maps to their corresponding branches of \(K_1\).
Hence its first norm, in the pro-three completion, is exactly
\[
 N_{3,1}w=\beta_1^{hk}U^z h_{11}^{\,hb}h_{59}^{\,hc}
\tag{A.174.13}
\]
for one \(z\in\Z_3\).
Indeed the displayed elements have precisely the same finite
divisor, and their quotient is a global unit of \(K_1\);
\(U\) generates its pro-three completion by A.165.
The possible sign is trivial in this completion.
There is no \(h_2\) factor: the norm divisor has only the
distinguished two-prime.

Apply A.174.11 to \(x=N_{3,2}w\).
It gives \(h(c-4b-3k)=0\) modulo nine.
Since \(h\) is invertible modulo nine, every defining relation
of \(G\) is killed by the rule in A.174.2.
That rule therefore defines the asserted homomorphism.
It is onto because it takes \(R_{59}\) to one.
The actual values of \(\chi_3\) in A.170 give
\(\chi_3(P)=0,\chi_3(R_{11})=2,\chi_3(R_{59})=1\);
these are exactly its reduction.

\emph{The literal principal relations and their field maps.}
Write \(\sigma_{11},\sigma_{59},\sigma_\Sigma\in K_2^\times\)
for the three integral elements of A.165.11.
In \(K_2\) their divisors are
\[
 P_{11,2,2},\qquad P_{59,2,2},\qquad
                       P_{11,0,2}+P_{59,0,2},
\]
respectively.  Each changed-prime branch of \(K_3/K_2\) is
totally ramified of degree three.  The actual embedding is
\(t=3d/v^3\), where \(v^{27}=3^8d^3\).
Thus their divisors after this embedding are three times the
corresponding displayed divisors at level three.
The rational \(q\) has divisor
\(9(P_{q,0,3}+P_{q,2,3})\).  Consequently the literal functions
\[
       f_q=\frac{\sigma_q^9}{q^2},\qquad
       w_0=\frac{\sigma_{11}\sigma_{59}}{\sigma_\Sigma^2}
\tag{A.174.14}
\]
satisfy
\[
 \operatorname{div}(f_q)=9\mathfrak R_q,\qquad
 \operatorname{div}(w_0)=3(\mathfrak R_{11}+\mathfrak R_{59}).
\tag{A.174.15}
\]
Divisors are written additively in this formula.
It proves \(9R_q=0\) and \(3(R_{11}+R_{59})=0\).
Their nonzero order-nine images under \(\theta\) prove that
each \(R_q\) has exact order nine.
Also \(3r=0\), with the explicit principal function
\(w_0/f_{59}^2\) for the divisor \(3r\).

Since \(\theta(R_{59})=1\) and \(R_{59}\) has order nine,
the map \(a\mapsto aR_{59}\) is an actual section of \(\theta\).
The equalities
\(R_{11}=5R_{59}+r\) and \(P=e'-3R_{59}\) give both direct
decompositions in A.174.4.  In particular
\(\ker(\theta|_{\mathcal R})=\langle r\rangle\), of order one
or three.

\emph{The extra two-prime relations and the exact balanced quotient.}
At two the extension \(K_3/K_2\) is unramified.
The level-two primes of absolute residue degrees \(1,2,6\)
extend respectively to the sum of the level-three primes of
degrees \(1,2\), the single prime of degree \(6\), and the single
prime of degree \(18\).
These are the actual root-tower factor maps of A.165--A.166.
Equivalently, over the old linear factor the cubic equation has
one linear and one unramified quadratic factor; on the other
two factors the three-primary residue unit has the maximal
available order and the next root gives an unramified cubic.
The three literal principal functions from level two therefore give
\[
 \begin{array}{c|c}
 \text{function, embedded in }K_3&\text{its divisor}\\ \hline
 \beta_2&P_{2,0,3}+P_{2,2,3}\\
 \beta_1/\beta_2&P_{2,6,3}\\
 2/(\alpha-1)&P_{2,18,3}.
 \end{array}
\tag{A.174.16}
\]
Thus the three balanced two-prime relations have classes
\(-3P,-6P,-18P\), and generate exactly \(\langle3P\rangle\).
There is no nonzero balanced relation at the unique three-prime.
The complete balanced relation subgroup of \(A\) is consequently
\[
              \mathcal R+\langle3P\rangle\subset G.
\tag{A.174.17}
\]
This proves the actual quotient and image formula in A.174.4:
after the split \(\Z/9\) factor is removed,
\(3P=3e'\), because \(9R_{59}=0\).

The quotient \(\langle e',r\rangle/\langle r\rangle\) is a
finite cyclic three-group generated by the image of \(e'\).
Its generator is zero modulo three precisely when the group is
zero: multiplication by three is surjective on a finite
three-group only when that group is trivial.
Applied to A.174.4, this proves
\[
 e_3=0\quad\Longleftrightarrow\quad
 e'\in\langle r\rangle
 \quad\Longleftrightarrow\quad
 e'+jr=0\text{ for some }j=0,1,2.
\tag{A.174.18}
\]
The last three classes are exactly those of A.174.5.
The residue degrees \(1,2\) at the changed primes make each
\(\mathfrak R_q\) have norm one, proving the norm assertion.
Vanishing of the three-primary class of a fractional ideal is
equivalent to a principal power of exponent prime to three;
this retains the difference between the integral principal ideal
test and the exact three-primary criterion.

For a smaller representative in the third test one may replace
the exponent \(3-5\cdot2=-7\) by \(2\).
The resulting ideal is \(\mathfrak I_2(f_{59})\), by
A.174.15, so the replacement has this specified principal
multiplier and the same class.
The finite affine congruences, all principal divisor vectors,
the split order-nine factor and the finite cyclic-group reduction
are checked in \path{code/two_inert_order_nine_balance_certificate.py}.
The recorded ideal searches are positive-only searches for these
three tests and do not add a vanishing assertion.
\end{proof}
